\chapter{Quantum Machine Learning}

\section{Mysteries of machine learning}


In recent decades the rapid development of cheap and powerful classical processing units has transformed all aspects of society. Artificial intelligence has been pondered by humanity from as far back as the 1800's but has remained the realm of science fiction until numerous great advances in the capability of computers has facilitated many implementations of AI across various industries from medicine to the military. Many of these AI technologies are based on machine learning (ML). 

\section{QC is a lot like machine learning}

There are quite a few similarities between quantum computing \& machine learning. They are both:

\begin{itemize}
    \item Built on linear algebra
    \item Held back by expensive computational resource
    \item Sometimes difficult to explain outcomes 
    \item Likely to go through period of hype \& stagnation as technological development pace changes
    \item Often associated with optimisation problems 
    \item Typically developed in a user-friendly environment using Python \& Jupyter notebooks 
\end{itemize}


It's an open question as to whether or not quantum computing can improve machine learning. There have been some early experiments showing quantum ML outperforming classical ML, but there have equally been experiments showing classical machine learning outperforming quantum machine learning. 

 Given the very limited number (\& quality) of qubits available there haven't been any large-scale tests yet. This might change in the near future...

\section{Why machine learning could benefit from a quantum advantage}

There are a few reasons motivating the use of quantum computers for machine learning:
\begin{itemize}
    \item Better sampling from true randomness of quantum states
    \item Higher dimensionality from $2^n$ dimension only requiring $n$ qubits
    \item Entanglement could be better for modelling correlations in joint-probability distributions
    \item For data generated by quantum systems can be used much \href{https://www.science.org/doi/10.1126/science.abn7293}{more effectively on a quantum model}
\end{itemize}

\section{Variants of QML} 

\subsection{Variational Quantum Eigensolver}

The variational quantum eigensolver (VQE) is one of the most well-known quantum algorithms. 

\subsection{QAOA}

\subsection{Generative Quantum Modelling}

\subsection{Quantum SVM} 

\section{Additional considerations}

There are quite a few things to consider with QML 
\begin{itemize}
    \item There are no proven advantages to using QML over classical ML 
    \item Quantum computers struggle to load large classical data sets 
    \item Many proposed quantum algorithms depend on quantum ram, which has yet to be scaled and integrated into QPU's. 
    \item There is significant academic scepticism, \href{https://www.youtube.com/watch?v=5UsJV2BNj2U&list=PLOFEBzvs-VvppIb0jg5_aDbmFs36DXD9w}{even in companies that promote QML}
    \item If a classical ML solution works perfectly well for a given problem, it is much better to use that than a quantum version. 
\end{itemize}


Some optimism: 

Quantum machine learning is still a very new field. It may take quite some time to offer benefits, but that doesn't mean it isn't worth exploring.

Conventional ML doesn't have much theoretical guarantees, but we know it works really well in practice. Given there are so many diverse applications of ML, one of them is likely to be well-suited to a quantum model.

Even if our current QML approaches don't offer any value at the moment, further research may uncover models better suited to the hardware we have now...

\subsection{Additional resources}
\begin{itemize}
    \item \href{https://github.com/Christophe-pere/Roadmap-to-QML}{GitHub repository listing papers, books and online courses}
    \item \href{https://www.youtube.com/watch?v=5UsJV2BNj2U&list=PLOFEBzvs-VvppIb0jg5_aDbmFs36DXD9w}{A panel discussion from 2021 on the Future of Quantum Machine Learning}
\end{itemize}
